\documentclass[11pt]{article}
\usepackage{amsmath,amssymb,amsthm}
\usepackage[margin=1.1in]{geometry}
\newtheorem{proposition}{Proposition}
\newtheorem{remark}{Remark}

\title{The Multivariate Gaussian}
\author{Yuval Lavie}
\date{\today}

\begin{document}
\maketitle

\noindent
$x \sim \mathcal{N}(\mu, \Sigma)$ on $\mathbb{R}^d$:
\begin{equation}
  f(x) = \frac{1}{(2\pi)^{d/2} |\Sigma|^{1/2}}
  \exp\!\Bigl(-\tfrac12 (x-\mu)^{\!\top} \Sigma^{-1} (x-\mu)\Bigr),
  \label{eq:density}
\end{equation}
$\Sigma$ symmetric positive definite. The level sets are ellipsoids:
$\Sigma$'s eigenvectors are their axes, its eigenvalues the variances
along them (drawn by the companion script as the $1\sigma/2\sigma$
ellipses over a sample).

\section{Closure under linear maps}

\begin{proposition}
If $x \sim \mathcal{N}(\mu, \Sigma)$ then
$Ax + b \sim \mathcal{N}(A\mu + b,\, A \Sigma A^{\!\top})$.
\end{proposition}
This is the property that makes Gaussians ``the linear algebra of
probability'': marginals, sums, and projections of Gaussians stay
Gaussian, with moments you can compute by matrix arithmetic. Verified
by sampling with a random $A$.

\section{Conditioning is linear regression}

Partition $x = (x_0, x_1)$ with blocks of $\Sigma$. Then
\begin{equation}
  \boxed{\;x_0 \mid x_1 \;\sim\; \mathcal{N}\bigl(
  \mu_0 + \Sigma_{01}\Sigma_{11}^{-1}(x_1 - \mu_1),\;
  \Sigma_{00} - \Sigma_{01}\Sigma_{11}^{-1}\Sigma_{10}\bigr)\;}
  \label{eq:cond}
\end{equation}
--- the conditional mean is \emph{linear} in $x_1$, and the
conditional covariance does not depend on $x_1$ at all
(homoscedasticity). Verified by binning a $2\times10^5$ sample: the
binned conditional means lie on the predicted line to $0.05$, the
conditional variance is flat at the predicted value.

\begin{remark}[Why regression is linear]
\eqref{eq:cond} is the generative justification of
Algorithms/regression/linear: if $(x, y)$ are jointly Gaussian, the
best predictor of $y$ from $x$ IS a line, and least squares estimates
exactly $\Sigma_{01}\Sigma_{11}^{-1}$.
\end{remark}

\section{Maximum likelihood}

For $n$ i.i.d.\ samples, maximizing \eqref{eq:density}'s
log-likelihood gives
\begin{equation}
  \boxed{\;\hat\mu = \bar x, \qquad
  \hat\Sigma = \frac1n \sum_i (x_i - \bar x)(x_i - \bar x)^{\!\top}\;}
  \label{eq:mle}
\end{equation}
The $1/n$ makes $\hat\Sigma$ biased:
$\mathbb{E}[\hat\Sigma] = \frac{n-1}{n}\Sigma$ (the sample mean has
eaten one degree of freedom) --- the 1D gaussian folder's lesson as a
matrix statement, measured at $n = 5$: empirical factor $0.80$ against
$(n-1)/n = 0.8$.

\end{document}
